\documentclass[a4paper,12pt]{extarticle}
\usepackage[T2A]{fontenc}
\usepackage[utf8]{inputenc}
\usepackage[russian]{babel}
\usepackage{geometry}
\geometry{margin=22mm}
\pagestyle{empty}
\begin{document}
\section*{Лекция и лабораторная работа № 4 · ключ преподавателя}
Студент выполняет вопросы 1–6 и вопросы 7–10 только по своему варианту
лабораторной. За верный ответ — 1 балл; максимум 10 баллов. Тест не является
допуском к защите лабораторной работы.

\subsection*{Общий блок}
\begin{center}
\begin{tabular}{|c|c|l|}\hline
Вопрос & Ответ & Проверяемое понятие \\ \hline
1 & Б & CPU выполняет инструкции \\ \hline
2 & А & Код программы во время исполнения находится в RAM \\ \hline
3 & А & Адрес и значение ячейки — разные вещи \\ \hline
4 & Б & Регистр — рабочая память внутри CPU \\ \hline
5 & Б & Элементы массива расположены подряд \\ \hline
6 & А & Стек хранит адрес возврата и данные вызова \\ \hline
\end{tabular}
\end{center}

\subsection*{Вариант 1 · C, GCC и GDB}
\begin{center}
\begin{tabular}{|c|c|l|}\hline
7 & Б & Изменённый исходник нужно пересобрать \\ \hline
8 & А & \texttt{-g} добавляет отладочную информацию \\ \hline
9 & А & GDB позволяет наблюдать состояние программы \\ \hline
10 & Б & \texttt{0! = 1} \\ \hline
\end{tabular}
\end{center}

\subsection*{Вариант 2 · NASM, BIOS и QEMU}
\begin{center}
\begin{tabular}{|c|c|l|}\hline
7 & В & Загрузочный сектор занимает 512 байтов \\ \hline
8 & Б & В текстовом VGA хранятся символ и атрибут \\ \hline
9 & Б & Программный вызов BIOS и аппаратный IRQ \\ \hline
10 & В & Без записи на диск оценка остаётся только в RAM \\ \hline
\end{tabular}
\end{center}
\end{document}
